\documentclass[11pt,a4paper]{article}
\usepackage[T1]{fontenc}
\usepackage[utf8]{inputenc}
\usepackage{lmodern,amsmath,amssymb}
\usepackage[margin=25mm]{geometry}
\usepackage[hidelinks]{hyperref}
\hypersetup{pdftitle={The ground-state measure is the time-slice marginal of the Euclidean measure, and that imp},pdfauthor={navstokgap research working notes}}
\newcommand{\tightlist}{\setlength{\itemsep}{0pt}\setlength{\parskip}{0pt}}
\setlength{\emergencystretch}{3em}
\setcounter{secnumdepth}{0}
\begin{document}
\hypertarget{the-ground-state-measure-is-the-time-slice-marginal-of-the-euclidean-measure-and-that-imports-the-local-large-field-estimate}{%
\section{The ground-state measure is the time-slice marginal of the
Euclidean measure, and that imports the local large-field
estimate}\label{the-ground-state-measure-is-the-time-slice-marginal-of-the-euclidean-measure-and-that-imports-the-local-large-field-estimate}}

The obstruction recorded in \href{magnetic-energy-identities.md}{the
identities note} is removed by an identification rather than by an
estimate. For the Kogut--Susskind Hamiltonian obtained as the
\(a_t\to0\) limit of the anisotropic Wilson transfer matrix,
\[\big|\Omega(U)\big|^2\,d\mu(U)\ =\ \lim_{T\to\infty}\
\text{marginal at time }0\text{ of }\ \frac{e^{-S_E[\mathcal U]/\hbar}\,\prod d\mu}{Z_T},\]
so the ground-state measure of the Hamiltonian theory is the
distribution of one time slice under the four-dimensional Euclidean
Wilson measure. Every estimate proved for that measure therefore holds
for \(|\Omega|^2\) verbatim, including the local ones that the
eigenvalue equation cannot reach, because the Euclidean weight
\(e^{-S_E/\hbar}=\prod_Pe^{-\beta_P/\hbar}\) factorizes over
four-dimensional plaquettes while the eigenvalue equation carries the
extensive \(E_0\). The estimate the route needs is then available in its
standard form: a configuration whose field strength averaged over a
block of side \(\ell\) reaches \(\eta/\ell^2\) costs Euclidean action
\[\frac{S_E}{\hbar}\ \simeq\ \frac{1}{4g^2}\!\int_{\text{block}}\!\big|F\big|^2\,d^4x
\ \simeq\ \frac{\eta^2}{4g^2},\] \textbf{independent of \(\ell\) and of
the lattice spacing}, so the suppression is \(e^{-c\eta^2/g(\ell)^2}\)
at every scale, with the running coupling. This matches, with the same
exponent structure, the free-field Gaussian tail
\(\exp[-8\pi^2\eta^2/g^2]\) computed in
\href{typical-field-strength-window.md}{the typical-field note} §3,
which is a check on both. With the identification, the decimation step
of the route controls its truncation error off a set of measure
\(e^{-c\eta^2/g^2}\), and what remains is the treatment of that set: the
large-field problem in its usual form, now reached honestly rather than
assumed. Constants explicit; the identification is standard and cited;
nothing promoted.

\hypertarget{the-transfer-matrix-and-the-limit}{%
\subsection{1. The transfer matrix and the
limit}\label{the-transfer-matrix-and-the-limit}}

On an anisotropic lattice with spatial spacing \(a\) and temporal
spacing \(a_t\), the Wilson action separates into temporal and spatial
plaquettes,
\[\frac{S_E}{\hbar}=\frac{1}{g^2}\Big[\frac{a}{a_t}\sum_{P\ \rm temporal}
\big(N-\operatorname{Re}\operatorname{tr}U_P\big)
+\frac{a_t}{a}\sum_{P\ \rm spatial}
\big(N-\operatorname{Re}\operatorname{tr}U_P\big)\Big]\cdot\frac{2}{1},\]
in the normalization of \href{mass-gap-obligations-lattice.md}{the
obligations map} with \(g\) dimensionless. The transfer matrix
\(\mathcal T\) of this action is self-adjoint and strictly positive
(Lüscher, Commun. Math. Phys. 54 (1977) 283; Creutz, Phys. Rev.~D 15
(1977) 1128; both metadata level), and as \(a_t\to0\)
\[\mathcal T=\exp\Big[-\frac{a_t}{\hbar}\,H_{\rm KS}\Big]\big(1+O(a_t)\big),
\qquad
H_{\rm KS}=\frac{\hbar c}{a}\Big[\frac{g^2}2\sum_\ell(-\Delta_\ell)
+\frac2{g^2}\sum_p\big(N-\operatorname{Re}\operatorname{tr}U_p\big)\Big],\]
the temporal plaquettes producing the electric term through the
character expansion and the spatial ones passing directly to the
magnetic term.

\textbf{Proposition 1.} Let \(\Omega\) be the Perron vector of
\(\mathcal T\), which by T1 of
\href{mass-gap-obligations-lattice.md}{the obligations map} is the
positive ground state of \(H_{\rm KS}\) in the limit. Then for any
bounded \(F\) of a single time slice,
\[\big\langle\Omega,F\,\Omega\big\rangle
=\lim_{T\to\infty}\frac{\big\langle\chi,\mathcal T^{T}F\,\mathcal T^{T}\chi\big\rangle}
{\big\langle\chi,\mathcal T^{2T}\chi\big\rangle}
=\lim_{T\to\infty}\big\langle F\big\rangle_{S_E,\,[-T,T]},\] for any
\(\chi\) with \(\langle\chi,\Omega\rangle\neq0\); that is,
\(|\Omega|^2d\mu\) is the time-slice marginal of the Euclidean measure.

\emph{Proof.} Spectral decomposition of \(\mathcal T\) with the gap of
T1 gives the first equality; the middle expression is by construction
the Euclidean expectation of \(F\) inserted at time \(0\) in a slab of
height \(2T\) with boundary state \(\chi\), which is the second.
\(\square\)

Nothing here is new; the point is what it licenses.

\hypertarget{what-transfers}{%
\subsection{2. What transfers}\label{what-transfers}}

Any statement of the form ``the Euclidean measure assigns probability at
most \(p\) to the set of configurations whose restriction to a time
slice has property \(\mathcal A\)'' is, by Proposition 1, a statement
about \(|\Omega|^2\). In particular the local statements do, and the
reason the Hamiltonian formulation could not produce them itself is now
visible as a feature of the representation rather than of the physics:
the eigenvalue equation \(-A'\Delta\Omega+(B'V-e_0)\Omega=0\) carries
the extensive \(e_0\) and compares totals, while \(e^{-S_E/\hbar}\) is a
product over plaquettes and compares locally. Both describe the same
measure.

This closes the gap left open in
\href{magnetic-energy-identities.md}{the identities note} §5 and in
\href{typical-field-strength-window.md}{the typical-field note} §3, in
the only way those notes left available.

\hypertarget{the-large-field-estimate-at-a-block-scale}{%
\subsection{3. The large-field estimate at a block
scale}\label{the-large-field-estimate-at-a-block-scale}}

\textbf{The action cost is scale-invariant.} Let a configuration have
field strength of magnitude \(\eta/\ell^2\) coherently over a
four-dimensional block of side \(\ell\). Its Euclidean action is
\[\frac{S_E}{\hbar}=\frac1{4g^2}\int\big|F^a_{\mu\nu}\big|^2d^4x
\ \simeq\ \frac{1}{4g^2}\Big(\frac{\eta}{\ell^2}\Big)^2\ell^4=\frac{\eta^2}{4g^2},\]
independent of \(\ell\) and of \(a\). This is the reason large-field
estimates take the same form at every scale, with only the coupling
running: the suppression is
\[\mathbb P\Big(\big|F\big|_{\rm block}\ \ge\ \frac{\eta}{\ell^2}\Big)
\ \lesssim\ \exp\Big[-\frac{c\,\eta^2}{g(\ell)^2}\Big],\] with \(c\) a
pure number and \(g(\ell)\) the coupling at the block scale.

\textbf{Consistency check.} The free-field computation of
\href{typical-field-strength-window.md}{the typical-field note} §3 gave,
for the flowed magnetic field at radius \(\sqrt{8t}=\ell\),
\[\mathbb P\Big(|b_t|>\frac{\eta}{t}\Big)\simeq\exp\Big[-\frac{8\pi^2\eta^2}{g^2}\Big],\]
the same \(\eta^2/g^2\) in the exponent, obtained by an independent
route from the Gaussian ground state. The two agree in structure and fix
\(c=8\pi^2\) in the free limit.

\textbf{Why the lattice-scale version is empty.} At \(\ell=a\) the
plaquette angles are of order one, the action per plaquette is bounded
by \(4N/g^2\), and the Chebyshev bound
\(\mathbb P(S_R\ge s)\le\exp[-s+4N|R|/g^2]\) is vacuous, as noted in
\href{typical-field-strength-window.md}{the typical-field note}. The
estimate has content only for the block-averaged field at \(\ell\gg a\),
where the individual plaquette angles are small, of order
\(a^2\eta/\ell^2\), and the cost accumulates over \((\ell/a)^4\) of them
to the scale-invariant total above. The large-field condition is a
statement about coherence across a block, not about any single
plaquette.

\hypertarget{where-this-leaves-the-route}{%
\subsection{4. Where this leaves the
route}\label{where-this-leaves-the-route}}

With Proposition 1 and Section 3, the decimation step of
\href{mass-gap-position.md}{the position note} §6 has its error
controlled outside a set of ground-state measure
\(e^{-c\eta^2/g(\ell)^2}\), which at weak coupling is small, and inside
that set the flow-truncation bound of
\href{flow-jacobian-truncation-error.md}{the Jacobian note} degrades
like \(e^{2t\|G\|_\infty}\). The two exponents compete:
\[\text{gain }e^{-c\eta^2/g^2}\qquad\text{against}\qquad\text{loss }e^{2\eta},\]
since \(t\|G\|_\infty\simeq\eta\) when the field reaches \(\eta/\ell^2\)
at the flow radius. The gain wins for \[\eta\ \gtrsim\ \frac{2g^2}{c},\]
so the large-field region can be defined as
\(\{\,|F|_{\rm block}\ge\eta_*/\ell^2\,\}\) with \(\eta_*=2g^2/c\), on
which the measure is at most \(e^{-4g^2/c}\) and outside which the
truncation is accurate. \textbf{At weak coupling \(\eta_*\to0\), so the
good region is almost everything}; the residual set carries measure
\(e^{-4g^2/c}\), which tends to one as \(g\to0\), and there the estimate
fails to be useful.

That last sentence is the honest statement of where the difficulty now
sits: the competition is between an exponential gain in \(\eta^2/g^2\)
and an exponential loss in \(\eta\), and at small \(g\) the crossover
\(\eta_*\) is small, so the excluded set is defined by a weak condition
and its measure bound \(e^{-4g^2/c}\) is close to one. Improving it
requires a sharper treatment of the truncation inside the large-field
region, which is the large-field expansion of the constructive
programme, or a truncation whose error grows more slowly than
\(e^{2t\|G\|_\infty}\).

\hypertarget{consequence-for-state}{%
\subsection{5. Consequence for STATE}\label{consequence-for-state}}

The identification of \(|\Omega|^2\) with the Euclidean time-slice
marginal is recorded, with the transfer-matrix references, and it
imports the local estimates that the Hamiltonian eigenvalue equation
could not produce. The large-field estimate itself has the
scale-invariant form \(e^{-c\eta^2/g^2}\), confirmed against the
independent free-field computation. The route's decimation step now has
a quantitative structure: good region with accurate truncation, bad
region of measure \(e^{-c\eta_*^2/g^2}\), and a crossover \(\eta_*\)
fixed by the competition between the two exponentials. The next question
is the first one that is genuinely about the large-field region rather
than around it: can the flow-truncation error inside it be bounded by
something weaker than \(e^{2t\|G\|_\infty}\), for instance by using that
the flow contracts the action monotonically even there?
\end{document}
